% Theme: Signature  --  the current @quantum_danesh look, in Palatino
% Build:  ./build.sh example-signature.tex   (or compile on Overleaf)
\documentclass[12pt]{beamer}
\usetheme{igsignature}

\ighandle{@quantum_danesh}
\igseries{Qubits desde cero}
\igissue{02}

% A Stern-Gerlach magnet: \pic at (x,y) {sg={axis/arrow angle/arrow colour}};
\tikzset{pics/sg/.style args={#1/#2/#3}{code={
  \draw[igAccentB, line width=0.7pt, fill=igPanel, rounded corners=1.6mm] (-7,-4.6) rectangle (7,4.6);
  \draw[igMuted!60, line width=0.4pt] (-3.4,0) circle (2.6);
  \draw[#3, line width=0.9pt, -{Latex[length=1.3mm]}] (-3.4,0) ++(#2+180:1.6) -- ++(#2:3.4);
  \node[text=igFG, font=\footnotesize] at (3.2,2.1) {$+#1$};
  \node[text=igFG, font=\footnotesize] at (3.2,-2.1) {$-#1$};
  \node[text=igMuted, font=\small] at (0,-7.5) {SG$_{#1}$};}}}
\newcommand\sgpass[2]{% from the + output of one magnet to the next
  \draw[igAccentB, line width=0.8pt] (#1+7,2.1) .. controls +(4,0) and +(-4,0) .. (#2-7,0);}
\newcommand\sgblock[1]{% the - output of a magnet, blocked
  \draw[igMuted, line width=0.8pt] (#1+7,-2.1) -- (#1+10,-2.1);
  \draw[igMuted, line width=1.6pt] (#1+10,-3.6) -- (#1+10,-0.6);}
\newcommand\sgout[3]{% both outputs of the last magnet, with their probabilities
  \draw[igAccentB, line width=0.8pt] (#1+7,2.1) -- (#1+10,4.5);
  \draw[igAccentB, line width=0.8pt] (#1+7,-2.1) -- (#1+10,-4.5);
  \node[anchor=west, fill=igHL, draw=igTrack, rounded corners=1.2mm, text=igAccent,
    font=\small, inner sep=1mm] at (#1+10,4.5) {#2};
  \node[anchor=west, fill=igHL, draw=igTrack, rounded corners=1.2mm, text=igAccent,
    font=\small, inner sep=1mm] at (#1+10,-4.5) {#3};}

\begin{document}

\coverslide{Qubits desde cero}
  {Medir cambia el estado}
  {Tres imanes de Stern--Gerlach, uno detrás de otro.}

\begin{frame}{Un imán, dos salidas}{El experimento}
\begin{center}
\begin{tikzpicture}[x=1mm,y=1mm,scale=1.45,transform shape]
  \draw[igAccentB, line width=0.8pt] (0,0) -- (21,0);
  \pic at (28,0) {sg={z/90/igAccentB}};
  \sgout{28}{50\,\%}{50\,\%}
\end{tikzpicture}
\end{center}
Un haz de átomos de plata atraviesa un imán no uniforme y se separa en
\hl{solo dos} manchas: $+z$ y $-z$.
\[
  \ket{\psi} = \alpha\ket{+z} + \beta\ket{-z},
  \qquad P(\pm z) = |\alpha|^2,\ |\beta|^2
\]
\end{frame}

\begin{frame}{Filtramos $+z$}{Primer paso}
\begin{center}
\begin{tikzpicture}[x=1mm,y=1mm,scale=1.3,transform shape]
  \draw[igAccentB, line width=0.8pt] (0,0) -- (8,0);
  \pic at (15,0) {sg={z/90/igAccentB}};
  \sgblock{15}
  \sgpass{15}{43}
  \pic at (43,0) {sg={z/90/igAccentB}};
  \sgout{43}{100\,\%}{0\,\%}
\end{tikzpicture}
\end{center}
Si bloqueamos $-z$ y volvemos a medir $z$, \hl{siempre} sale $+z$.
La primera medida ha preparado el estado.
\end{frame}

\begin{frame}{Ahora medimos $x$}{Segundo paso}
\begin{center}
\begin{tikzpicture}[x=1mm,y=1mm,scale=1.3,transform shape]
  \draw[igAccentB, line width=0.8pt] (0,0) -- (8,0);
  \pic at (15,0) {sg={z/90/igAccentB}};
  \sgblock{15}
  \sgpass{15}{43}
  \pic at (43,0) {sg={x/0/igAccent}};
  \sgout{43}{50\,\%}{50\,\%}
\end{tikzpicture}
\end{center}
\[
  \ket{+z} = \frac{1}{\sqrt2}\big(\ket{+x} + \ket{-x}\big)
\]
Cada resultado de $x$ sale con probabilidad \hl{$\tfrac12$}.
\end{frame}

\begin{frame}{Volvemos a medir $z$}{La sorpresa}
\begin{center}
\begin{tikzpicture}[x=1mm,y=1mm,scale=1.15,transform shape]
  \draw[igAccentB, line width=0.8pt] (-2,0) -- (0,0);
  \pic at (7,0) {sg={z/90/igAccentB}};
  \sgblock{7}
  \sgpass{7}{31}
  \pic at (31,0) {sg={x/210/igAccent}};
  \sgblock{31}
  \sgpass{31}{55}
  \pic at (55,0) {sg={z/90/igAccentB}};
  \sgout{55}{50\,\%}{50\,\%}
\end{tikzpicture}
\end{center}
Aunque antes seleccionamos $+z$, después de medir y seleccionar $+x$, una
nueva medición en $z$ puede dar \hl{cualquiera de los dos resultados}.
\end{frame}

\begin{frame}{Medir no es mirar}{La moraleja}
\keyeq{\big|\braket{\pm z}{+x}\big|^2 = \mhl{\tfrac12}}
Medir $x$ borra lo que sabíamos de $z$, porque los dos observables no
conmutan:
\[
  [S_x, S_z] = -i\hbar\, S_y \neq 0 .
\]
\end{frame}

\closingslide{Medir es intervenir}
  {En el próximo post: qué es exactamente un qubit.}

\end{document}
